% =====================================================================
\chapter{Modulo 12 -- Pile e Code}
% =====================================================================

% ---------------------------------------------------------------------
\section{Teoria \& Analogia Chiave}
% ---------------------------------------------------------------------

Pile e code non sono strutture ``nuove'': sono \textbf{liste (o array) a cui si
vietano quasi tutte le operazioni}. Si può entrare e uscire solo da punti precisi.
Proprio questa limitazione le rende utilissime: garantisce un ordine di uscita, e
rende ogni operazione $O(1)$.

\subsection*{La pila (stack): LIFO}

\begin{analogia}[: la pila di piatti]
In una pila di piatti appoggi un piatto \textbf{sopra} agli altri e, quando ne prendi
uno, prendi quello \textbf{in cima}: l'ultimo appoggiato. Per prendere quello in fondo
devi prima togliere tutti gli altri. Questa regola si chiama \textbf{LIFO}: \emph{Last
In, First Out}, l'ultimo entrato è il primo a uscire.
\end{analogia}

\begin{definizione}[: pila]
Una \textbf{pila} è una sequenza a cui si accede da \textbf{un solo estremo}, la
\textbf{cima}. Operazioni:
\begin{itemize}
  \item \cc|push(x)|: inserisce \cc|x| in cima;
  \item \cc|pop()|: rimuove (e di solito restituisce) l'elemento in cima;
  \item \cc|top()|: legge l'elemento in cima senza toglierlo;
  \item \cc|vuota()|, e \cc|piena()| se la capacità è limitata.
\end{itemize}
Fare \cc|pop| o \cc|top| su una pila vuota è un errore chiamato \textbf{underflow};
fare \cc|push| su una pila piena è un \textbf{overflow}.
\end{definizione}

La pila la conosci già: lo \textbf{stack} delle chiamate di funzione (Moduli 3 e 8) è
proprio una pila di record di attivazione. Altri usi tipici: annulla/ripeti negli
editor, il pulsante ``indietro'' del browser, il controllo delle parentesi, la
valutazione delle espressioni, invertire qualcosa.

\textbf{Pila su array} (come nelle slide). Un array \cc|dati| e un indice \cc|cima| che
parte da $-1$ (pila vuota):
\begin{itemize}
  \item \cc|push|: verifica che non sia piena, incrementa \cc|cima|, scrivi
        \cc|dati[cima]|;
  \item \cc|pop|: verifica che non sia vuota, leggi \cc|dati[cima]|, decrementa \cc|cima|.
\end{itemize}
Il fondo della pila è la posizione 0. La capacità è fissata (o si raddoppia come negli
Esercizi~4.10 e 5.10).

\textbf{Pila su lista concatenata.} Basta una lista semplice in cui la \textbf{cima è
la testa}: \cc|push| è un inserimento in testa, \cc|pop| una rimozione in testa. Entrambi
$O(1)$, nessun limite di capacità. La pila è vuota quando la testa è \cc|nullptr|.
\begin{lstlisting}
void push(const T& x) { cima = new Nodo<T>{x, cima}; }
T pop() {
    if (vuota()) throw underflow_error("pop su pila vuota");
    Nodo<T>* vecchia = cima;
    T x = vecchia->valore;     // salvo il valore PRIMA del delete
    cima = cima->succ;
    delete vecchia;
    return x;
}
\end{lstlisting}

\subsection*{La coda (queue): FIFO}

\begin{analogia}[: la fila alla posta]
Chi arriva si mette in \textbf{fondo} alla fila; lo sportello serve chi è in
\textbf{testa}, cioè chi è arrivato per primo. Regola \textbf{FIFO}: \emph{First In,
First Out}, il primo entrato è il primo a uscire.
\end{analogia}

\begin{definizione}[: coda]
Una \textbf{coda} è una sequenza in cui si inserisce da un estremo (il fondo) e si
rimuove dall'altro (la testa). Operazioni delle slide:
\begin{itemize}
  \item \cc|enqueue(x)|: aggiunge \cc|x| in fondo;
  \item \cc|dequeue()|: rimuove l'elemento in testa;
  \item \cc|front()|: legge l'elemento in testa senza toglierlo;
  \item \cc|isEmpty()|, e \cc|isFull()| se la capacità è limitata.
\end{itemize}
\end{definizione}

Usi tipici: code di stampa, gestione delle richieste a un server, simulazioni di
sportelli, buffer tra chi produce dati e chi li consuma, visita in ampiezza di alberi e
grafi (Modulo 14).

\textbf{Coda su lista concatenata.} Servono \textbf{due} puntatori: \cc|testa| (da cui
si esce) e \cc|fondo| (da cui si entra). Con il solo \cc|testa|, \cc|enqueue| dovrebbe
scorrere tutta la lista: $O(n)$. Con \cc|fondo| entrambe le operazioni sono $O(1)$.
Attenzione ai casi limite: inserendo in una coda vuota il nuovo nodo diventa sia testa
sia fondo; togliendo l'ultimo elemento vanno azzerati entrambi.
\begin{lstlisting}
void enqueue(const T& x) {
    Nodo<T>* nuovo = new Nodo<T>{x, nullptr};
    if (vuota()) testa = nuovo;         // era vuota: è anche la testa
    else fondo->succ = nuovo;           // altrimenti lo aggancio dopo l'ultimo
    fondo = nuovo;
}
\end{lstlisting}

\textbf{Coda su array circolare.} Con un array, se ogni \cc|dequeue| spostasse avanti
tutti gli elementi costerebbe $O(n)$. Il trucco è non spostare niente: si tengono
l'indice del primo elemento (\cc|inizio|) e il numero di elementi (\cc|n|), e si fa
``girare'' l'array come un orologio usando il resto della divisione:
\begin{itemize}
  \item \cc|enqueue|: \cc|dati[(inizio + n) % N] = x; n++;|
  \item \cc|dequeue|: \cc|x = dati[inizio]; inizio = (inizio + 1) % N; n--;|
\end{itemize}
Quando si arriva in fondo all'array, \cc|% N| riporta all'indice 0, nelle celle già
liberate. Coda vuota: \cc|n == 0|; coda piena: \cc|n == N|.

\begin{analogia}[: il nastro del sushi]
L'array circolare è il nastro di un ristorante di sushi: i piatti non si spostano
all'indietro quando qualcuno ne prende uno; semplicemente il ``punto di inizio'' avanza,
e dopo l'ultimo posto del nastro c'è di nuovo il primo.
\end{analogia}

\subsection*{Riepilogo}

\begin{center}
\begin{tabular}{lll}
\toprule
 & \textbf{Pila} & \textbf{Coda} \\
\midrule
politica & LIFO & FIFO \\
inserimento & \cc|push| (in cima) & \cc|enqueue| (in fondo) \\
rimozione & \cc|pop| (dalla cima) & \cc|dequeue| (dalla testa) \\
lettura & \cc|top| & \cc|front| \\
su lista & testa = cima & puntatori a testa \emph{e} fondo \\
su array & indice della cima & array circolare (\cc|% N|) \\
complessità & tutto $O(1)$ & tutto $O(1)$ \\
\bottomrule
\end{tabular}
\end{center}

\subsection*{Due applicazioni classiche della pila}

\textbf{Parentesi bilanciate.} Scorri il testo: ogni parentesi aperta va sulla pila;
ogni parentesi chiusa deve corrispondere a quella in cima (che si toglie). Errore se
trovi una chiusa con la pila vuota o del tipo sbagliato, oppure se alla fine la pila non
è vuota. Funziona perché l'ultima parentesi aperta è la prima che va chiusa: LIFO.

\textbf{Espressioni in notazione postfissa} (``polacca inversa''): l'operatore si
scrive \emph{dopo} i suoi operandi. \texttt{(3 + 4) * 2} diventa \texttt{3 4 + 2 *}. Non
servono parentesi né regole di precedenza. Per calcolarla: i numeri vanno sulla pila;
quando arriva un operatore si tolgono due numeri, si calcola e si rimette il risultato.
Attenzione: il primo numero estratto è il \emph{secondo} operando
(\texttt{8 2 -} vale $8 - 2$, non $2 - 8$).

\begin{inparole}
\begin{itemize}
  \item Pila = piatti: LIFO, si lavora solo in cima (\cc|push|, \cc|pop|, \cc|top|).
  \item Coda = fila alla posta: FIFO, si entra in fondo e si esce dalla testa
        (\cc|enqueue|, \cc|dequeue|, \cc|front|).
  \item Su lista: pila = inserimento e rimozione in testa; coda = puntatori a testa e fondo.
  \item Su array: pila con l'indice della cima; coda circolare con \cc|% N|.
  \item Togliere da una struttura vuota = underflow: lancia un'eccezione.
\end{itemize}
\end{inparole}

% ---------------------------------------------------------------------
\section{Rappresentazione Visiva della Memoria (TikZ)}
% ---------------------------------------------------------------------


eedspace{6cm}
\textbf{Pila e coda dopo l'arrivo di A, B, C, D} (il pattern).
\begin{center}
\begin{tikzpicture}
  % pila
  \foreach \v [count=\i] in {A,B,C,D}
    \node[cella,fill=heapcol,minimum width=1.6cm] (p\i) at (0,\i*0.75) {\v};
  \draw[freccia] (1.6,4.4) -- node[above,etichetta]{push} (0.9,3.4);
  \draw[freccia] (-0.9,3.4) -- node[above,etichetta]{pop} (-1.6,4.4);
  \node[etichetta,below=0.15cm of p1] {PILA: esce D (l'ultimo)};
  % coda
  \foreach \v [count=\i] in {A,B,C,D}
    \node[cella,fill=pubcol,minimum width=1.2cm] (c\i) at (4.6+\i*1.25,2) {\v};
  \node[etichetta,above=0.15cm of c1] {testa};
  \node[etichetta,above=0.15cm of c4] {fondo};
  \draw[freccia] ($(c1.west)+(0,0)$) -- node[above,etichetta]{dequeue} ++(-1.4,0);
  \draw[freccia] ($(c4.east)+(1.4,0)$) node[above left,etichetta]{enqueue} -- (c4.east);
  \node[etichetta,below=0.15cm of c2,xshift=0.6cm] {CODA: esce A (il primo)};
\end{tikzpicture}
\end{center}

\needspace{5cm}
\textbf{Coda su lista: i due puntatori.} Si entra da \cc|fondo|, si esce da \cc|testa|.
\begin{center}
\begin{tikzpicture}[nodoC/.style={draw,rectangle split,rectangle split horizontal,rectangle split parts=2,fill=heapcol,font=\ttfamily\small,minimum height=0.6cm}]
  \node[cella,fill=stackcol] (t) at (0,0.9) {testa};
  \node[cella,fill=stackcol] (f) at (7.2,0.9) {fondo};
  \node[nodoC] (a) at (1.6,-0.4) {A \nodepart{two} \phantom{x}};
  \node[nodoC] (b) at (3.6,-0.4) {B \nodepart{two} \phantom{x}};
  \node[nodoC] (c) at (5.6,-0.4) {C \nodepart{two} \phantom{x}};
  \node[font=\ttfamily\footnotesize] (z) at (7.4,-0.4) {nullptr};
  \draw[freccia] (t) -- (a.north);
  \draw[freccia] (f) -- (c.north);
  \draw[freccia] (a.two) -- (b.west);
  \draw[freccia] (b.two) -- (c.west);
  \draw[freccia] (c.two) -- (z.west);
\end{tikzpicture}
\end{center}

\needspace{5cm}
\textbf{Coda circolare con $N = 5$.} Dopo vari inserimenti e rimozioni, gli elementi
occupano le celle 3, 4, 0: il fondo ha ``fatto il giro''.
\begin{center}
\begin{tikzpicture}
  \foreach \v/\f [count=\i from 0] in {E/pubcol,{}/sfondo,{}/sfondo,C/pubcol,D/pubcol}
    \node[cellina,fill=\f,minimum width=1cm] (q\i) at (\i*1,0) {\v};
  \foreach \i in {0,...,4} \node[etichetta,below=0.05cm of q\i] {\i};
  \node[etichetta,above=0.15cm of q3] {inizio = 3};
  \node[etichetta,anchor=west,text width=6cm] at (5.4,0) {n = 3. Prossimo inserimento in \texttt{(3 + 3) \% 5 = 1}; prossima uscita: C.};
  \draw[freccia,dashed] (q4.north) to[out=60,in=120] (q0.north);
\end{tikzpicture}
\end{center}

\needspace{5cm}
\textbf{Valutare \texttt{3 4 + 2 *}}: lo stato della pila dopo ogni simbolo.
\begin{center}
\begin{tikzpicture}
  \foreach \tok/\stack [count=\i] in {3/{3},4/{3,4},+/{7},2/{7,2},*/{14}} {
    \node[font=\ttfamily\bfseries] at (\i*2.2,-0.9) {\tok};
    \foreach \v [count=\j] in \stack
      \node[cellina,fill=heapcol,minimum width=0.9cm] at (\i*2.2,\j*0.6-0.6) {\v};
  }
  \node[etichetta,anchor=east] at (1.3,-0.9) {simbolo letto};
  \node[etichetta,anchor=east] at (1.3,0.3) {pila (cima in alto)};
\end{tikzpicture}
\end{center}

% ---------------------------------------------------------------------
\section{Pattern di Codice di Laboratorio}
% ---------------------------------------------------------------------

Pila e coda generiche su nodi dinamici, con eccezioni di underflow e copia vietata
(\cc|= delete|). Lo stesso gruppo di persone entra nelle due strutture ed esce in
ordine opposto.

\codice{codice/m12/pattern.cpp}
\uscita{codice/m12/output/pattern.txt}

% ---------------------------------------------------------------------
\section{Tips \& Trappole da Esame}
% ---------------------------------------------------------------------

\begin{trappola}[: pop su struttura vuota]
Senza il controllo, \cc|pop()| su pila vuota dereferenzia \cc|nullptr|: crash. Controlla
sempre e lancia \cc|underflow_error| (o restituisci un \cc|bool|). Nel \cc|main|, cicla
con \cc|while (!p.vuota())|.
\end{trappola}

\begin{trappola}[: dimenticare il fondo nella coda]
Dopo l'ultimo \cc|dequeue|, \cc|testa| diventa \cc|nullptr| ma \cc|fondo| punta ancora
al nodo appena cancellato. Il successivo \cc|enqueue| funziona solo se controlla
\cc|vuota()| guardando \cc|testa|; altrimenti scrive in memoria liberata. Azzera
\cc|fondo| quando la coda si svuota.
\end{trappola}

\begin{trappola}[: leggere il valore dopo il delete]
In \cc|pop|/\cc|dequeue| il valore va copiato in una variabile \emph{prima} di
\cc|delete|, poi restituito.
\end{trappola}

\begin{trappola}[: ordine degli operandi nella postfissa]
Con \cc|b = p.pop(); a = p.pop();| il risultato di \texttt{-} e \texttt{/} è \cc|a - b|.
Invertirli dà risultati sbagliati solo per sottrazione e divisione, quindi il bug sfugge
facilmente ai test.
\end{trappola}

\begin{trappola}[: coda circolare piena o vuota?]
Con solo \cc|inizio| e \cc|fine| (senza contatore) le due situazioni ``piena'' e
``vuota'' hanno gli indici uguali. Usa un contatore \cc|n|, oppure lascia sempre una cella
vuota.
\end{trappola}

\begin{consiglio}[: riusare la lista]
All'esame puoi costruire pila e coda sopra la tua classe lista (pila = inserisci e togli
in testa; coda = inserisci in coda e togli in testa). Ma scriverle da zero con i nodi
richiede poche righe ed evita dipendenze: quasi sempre conviene.
\end{consiglio}

% ---------------------------------------------------------------------
\section{Quiz a Crocette (Prova Scritta)}
% ---------------------------------------------------------------------

\begin{quiz}{Si eseguono \texttt{push(1); push(2); push(3); pop(); push(4);} su una pila vuota. Cosa restituisce ora \texttt{top()}?}
  \item 1
  \item 2
  \item 3
  \item 4
\end{quiz}

\begin{quiz}{Si eseguono \texttt{enqueue(1); enqueue(2); enqueue(3); dequeue(); enqueue(4);} su una coda vuota. Cosa restituisce ora \texttt{front()}?}
  \item 1
  \item 2
  \item 3
  \item 4
\end{quiz}

\begin{quiz}{Quale politica segue una coda?}
  \item LIFO
  \item FIFO
  \item per priorità
  \item casuale
\end{quiz}

\begin{quiz}{Implementando una pila con una lista concatenata, dove conviene mettere la cima?}
  \item in testa alla lista
  \item in coda alla lista
  \item a metà lista
  \item indifferente: costa sempre $O(n)$
\end{quiz}

\begin{quiz}{Perché in una coda su lista concatenata si tiene anche un puntatore al fondo?}
  \item per poter scorrere la lista all'indietro
  \item per rendere \texttt{enqueue} $O(1)$ invece di $O(n)$
  \item per rendere \texttt{dequeue} $O(1)$
  \item non serve
\end{quiz}

\begin{quiz}{Coda circolare su array di dimensione 5, con \texttt{inizio = 3} e \texttt{n = 4}. In quale indice viene scritto il prossimo elemento?}
  \item 7
  \item 2
  \item 4
  \item 0
\end{quiz}

\begin{quiz}{Cosa succede facendo \texttt{pop} su una pila vuota?}
  \item restituisce 0
  \item si verifica un underflow
  \item si verifica un overflow
  \item la pila si riempie
\end{quiz}

\begin{quiz}{Quanto vale l'espressione postfissa \texttt{6 2 / 3 *}?}
  \item 1
  \item 4
  \item 9
  \item 36
\end{quiz}

\begin{quiz}{Quale struttura dati è più adatta per controllare se le parentesi di un'espressione sono bilanciate?}
  \item una coda
  \item una pila
  \item un array ordinato
  \item una lista circolare
\end{quiz}

\begin{quiz}{Qual è la complessità di \texttt{push}, \texttt{pop}, \texttt{enqueue} e \texttt{dequeue} nelle implementazioni viste?}
  \item $O(1)$
  \item $O(\log n)$
  \item $O(n)$
  \item $O(n)$ per la coda, $O(1)$ per la pila
\end{quiz}

\begin{soluzione}[: risposte commentate]
\begin{enumerate}[label=\textbf{\arabic*.}]
  \item \textbf{D}. Pila: 1, 2, 3; \cc|pop| toglie 3; poi entra 4, che sta in cima.
  \item \textbf{B}. Coda: 1, 2, 3; \cc|dequeue| toglie 1; in testa ora c'è 2.
  \item \textbf{B}.
  \item \textbf{A}. Inserimento e rimozione in testa costano $O(1)$; in coda la rimozione
        costerebbe $O(n)$.
  \item \textbf{B}.
  \item \textbf{B}. $(3 + 4) \% 5 = 2$.
  \item \textbf{B}.
  \item \textbf{C}. $6 / 2 = 3$, poi $3 \cdot 3 = 9$.
  \item \textbf{B}. L'ultima aperta deve essere la prima chiusa: LIFO.
  \item \textbf{A}. Con il puntatore al fondo anche la coda è tutta $O(1)$.
\end{enumerate}
\end{soluzione}

% ---------------------------------------------------------------------
\section{Laboratorio a Step (Simulazione d'Esame)}
% ---------------------------------------------------------------------

\begin{traccia}[: pile e code al lavoro]
Implementare da zero una pila e una coda generiche su nodi dinamici e una coda su array
circolare, e usarle per tre applicazioni: controllo delle parentesi, simulazione di uno
sportello e calcolatrice in notazione postfissa.
\end{traccia}

\begin{step}{la pila}
\cc|template <typename T> class Pila| su nodi, con \cc|push|, \cc|pop| (che lancia
\cc|underflow_error| se vuota), \cc|vuota| e distruttore.\\[0.3em]
\textbf{Concetto nuovo:} LIFO su lista.
\end{step}

\begin{step}{parentesi bilanciate}
Scrivi \cc|bool bilanciata(const string& s)| che controlla le parentesi
\texttt{() [] \{\}} usando una \cc|Pila<char>|. Provala su quattro espressioni, tra cui
\texttt{([)]}.\\[0.3em]
\textbf{Concetto nuovo:} usare la pila per ``ricordare'' in ordine inverso.
\end{step}

\begin{step}{la coda}
\cc|template <typename T> class Coda| con \cc|testa| e \cc|fondo|, \cc|enqueue|,
\cc|dequeue|, \cc|vuota|, \cc|size|.\\[0.3em]
\textbf{Concetto nuovo:} FIFO su lista con due puntatori.
\end{step}

\begin{step}{lo sportello}
Con \cc|srand(111222333)| simula 60 minuti di uno sportello: ogni minuto arriva un
cliente con probabilità 60\% (si mette in coda il suo minuto di arrivo); se lo sportello
è libero serve il primo della fila, con un tempo di servizio casuale tra 1 e 3 minuti.
Stampa arrivati, serviti, rimasti in fila, fila massima e attesa media.\\[0.3em]
\textbf{Concetto nuovo:} la coda come strumento di simulazione.
\end{step}

\begin{step}{coda circolare}
\cc|template <typename T, int N> class CodaCircolare| su array statico, con
\cc|isEmpty|, \cc|isFull|, \cc|enqueue| (lancia \cc|overflow_error| se piena) e
\cc|dequeue|. Mostra che dopo una rimozione l'elemento successivo riusa la cella 0.\\[0.3em]
\textbf{Concetto nuovo:} array circolare, operatore \cc|%|.
\end{step}

\begin{step}{calcolatrice postfissa}
Scrivi \cc|double valutaPostfissa(const string& espr)| che legge i simboli separati da
spazi con uno \cc|stringstream| e calcola il risultato con una \cc|Pila<double>|.
\cc|stod| converte una \cc|string| in \cc|double|.\\[0.3em]
\textbf{Concetto nuovo:} valutazione di espressioni con la pila.
\end{step}

\subsection*{Soluzione completa}
\codice{codice/m12/lab.cpp}
\uscita{codice/m12/output/lab.txt}

\begin{consiglio}[: cosa notare nella soluzione]
\begin{itemize}
  \item \cc|for (char c : s)| è il ``for per ogni elemento'': scorre tutti i caratteri
        della stringa senza indici.
  \item Nello sportello arrivano in media 0,6 clienti al minuto ma se ne servono circa
        0,5 (un servizio dura in media 2 minuti): la fila tende a crescere. Prova a
        cambiare la probabilità di arrivo e guarda come cambiano fila e attese.
  \item In \cc|valutaPostfissa|, se alla fine la pila contiene più di un numero
        l'espressione era malformata (mancavano operatori).
\end{itemize}
\end{consiglio}

% ---------------------------------------------------------------------
\section{Domande e Risposte Modello per l'Orale}
% ---------------------------------------------------------------------

\begin{orale}[D1. Cos'è una pila e come si implementa?]
Una pila è una struttura dati lineare a cui si accede da un solo estremo, la cima,
gestita con politica LIFO: l'ultimo elemento inserito è il primo a essere estratto. Le
operazioni sono \cc|push| (inserimento in cima), \cc|pop| (rimozione dalla cima), \cc|top|
(lettura della cima) e il test di pila vuota. Con un array si usa un indice che segna la
cima, incrementato a ogni \cc|push|; la capacità è fissa e la pila può essere piena. Con
una lista concatenata la cima è la testa della lista: \cc|push| e \cc|pop| sono
inserimento e rimozione in testa, in $O(1)$, senza limiti di capacità. Estrarre da una
pila vuota provoca un underflow.
\end{orale}

\begin{orale}[D2. Cos'è una coda e come si implementa?]
Una coda è una struttura lineare gestita con politica FIFO: gli elementi si inseriscono
in fondo (\cc|enqueue|) e si rimuovono dalla testa (\cc|dequeue|), nello stesso ordine in
cui sono entrati; \cc|front| legge l'elemento in testa. Con una lista concatenata si
mantengono due puntatori, alla testa e al fondo, così entrambe le operazioni sono
$O(1)$. Con un array si usa una coda circolare: un indice di inizio e il numero di
elementi, con gli indici calcolati modulo la dimensione, per non dover spostare elementi.
\end{orale}

\begin{orale}[D3. Pila su array o su lista: vantaggi e svantaggi?]
Su array l'accesso è molto veloce e non c'è spazio extra per i puntatori, ma la
capacità è fissata: la pila può riempirsi, oppure va riallocata raddoppiando, con un costo
occasionale $O(n)$. Su lista la dimensione è illimitata e ogni operazione è $O(1)$, ma
ogni elemento occupa memoria in più per il puntatore e ogni \cc|push| richiede
un'allocazione dinamica.
\end{orale}

\begin{orale}[D4. Fai un esempio di uso della pila.]
Il controllo delle parentesi bilanciate: si scorre l'espressione, si fa \cc|push| di ogni
parentesi aperta e, per ogni parentesi chiusa, si fa \cc|pop| verificando che corrisponda;
l'espressione è bilanciata se non si incontrano errori e alla fine la pila è vuota. Altri
esempi: la valutazione di espressioni in notazione postfissa, l'annulla/ripeti negli
editor e soprattutto lo stack delle chiamate di funzione, che è una pila di record di
attivazione.
\end{orale}

% ---------------------------------------------------------------------
\section{Esercizi al PC}
% ---------------------------------------------------------------------

\begin{esercizio}{Pila su array}{facile}
Implementa \cc|template <typename T, int CAP> class PilaArray| come nelle slide: indice
\cc|cima| che parte da $-1$, \cc|push| che verifica che la pila non sia piena,
\cc|pop| che verifica che non sia vuota. Provala con capacità 3 e quattro \cc|push|.
\end{esercizio}

\begin{esercizio}{Invertire con la pila}{facile}
Con una pila inverti una stringa carattere per carattere, e poi l'ordine delle parole di
una frase (leggi le parole con uno \cc|stringstream|).
\end{esercizio}

\begin{esercizio}{Coda di stampa}{facile}
Una coda di lavori di stampa (\cc|struct Lavoro| con nome del documento e pagine). La
stampante fa 20 pagine al minuto: stampa a che minuto inizia ogni lavoro.
\end{esercizio}

\begin{esercizio}{Cambio di base}{facile}
Converti un intero in base 2, 8 e 16 dividendo ripetutamente per la base e mettendo i
resti su una pila: escono nell'ordine giusto.
\end{esercizio}

\begin{esercizio}{Palindromi con pila e coda}{medio}
Inserisci le lettere di una frase sia in una pila sia in una coda (ignorando spazi,
punteggiatura e maiuscole). La frase è palindroma se pila e coda restituiscono sempre la
stessa lettera.
\end{esercizio}

\begin{esercizio}{Coda con due pile}{medio}
Implementa una coda usando solo due pile, \cc|entrata| e \cc|uscita|: \cc|enqueue| fa
\cc|push| su \cc|entrata|; \cc|dequeue|, se \cc|uscita| è vuota, vi travasa tutta
\cc|entrata| e poi fa \cc|pop|.
\end{esercizio}

\begin{esercizio}{Pila con minimo}{medio}
Una pila che, oltre a \cc|push| e \cc|pop|, restituisce in $O(1)$ il minimo degli elementi
presenti. Idea: ogni nodo memorizza anche il minimo della pila ``da lui in giù''.
\end{esercizio}

\begin{esercizio}{Annulla e ripeti}{medio}
Una classe \cc|Editor| con un testo e due pile di versioni precedenti: \cc|scrivi(s)|
salva la versione attuale nella pila \cc|annulla| e svuota \cc|ripeti|; \cc|undo()| e
\cc|redo()| spostano le versioni tra le due pile.
\end{esercizio}

\begin{esercizio}{Da infissa a postfissa}{difficile}
Converti espressioni infisse con cifre singole, \texttt{+ - * /} e parentesi in
postfisse: le cifre vanno subito in uscita; un operatore, prima di andare sulla pila,
fa uscire quelli in cima con priorità maggiore o uguale; una \texttt{)} fa uscire tutto
fino alla \texttt{(}. Poi valuta la postfissa.
\end{esercizio}

\begin{esercizio}{Pronto soccorso}{difficile}
Una coda con priorità: i pazienti hanno un codice (1 rosso, 2 giallo, 3 verde). Escono
prima i codici più urgenti e, a parità di codice, nell'ordine di arrivo. Realizzala con
una lista mantenuta ordinata all'inserimento.
\end{esercizio}

\subsection*{Soluzioni}
Le soluzioni commentate, con il codice e l'output reale, sono nel PDF separato
\textbf{Soluzioni} (\texttt{soluzioni-light.pdf} / \texttt{soluzioni-dark.pdf}),
allo stesso numero di esercizio. Il codice è anche in \texttt{codice/m12/}.
